% Generado por sistema/backend ReporteTex.java -- no editar a mano.
% Datos: CSV importado (titulados) / CSV importado (empleadores)
\subsubsection*{C11. Formación complementaria $\times$ interés en posgrado}

\textit{Paso 1. Describir el problema.} Establecer si haber cursado formación complementaria influye en el interés por seguir estudiando un posgrado.

\begin{itemize}[nosep]
  \item Variable independiente $X$: Formación complementaria.
  \item Variable dependiente $Y$: Interés en posgrado.
  \item Población: todos los titulados ($N = 148$). Criterio ARCU-SUR: 2.4.1. Nivel de significancia $\alpha = 0.05$.
  \item $H_0$: Las variables «Formación complementaria» (X) y «Interés en posgrado» (Y) son independientes.
  \item $H_1$: Existe dependencia significativa entre las variables «Formación complementaria» (X) y «Interés en posgrado» (Y).
\end{itemize}

\textit{Datos.} Primeros 10 de los 148 titulados que respondieron ambas preguntas (sin datos personales).

\begin{mitabla}{C11: datos de los primeros titulados}
\begin{tabular}{lll}
\toprule
Titulado & $X$: Formación complementaria & $Y$: Interés en posgrado \\
\midrule
T1 & Sí & Sí \\
T2 & Sí & Sí \\
T3 & Sí & Sí \\
T4 & Sí & Sí \\
T5 & Sí & Sí \\
T6 & Sí & Sí \\
T7 & Sí & Sí \\
T8 & Sí & Sí \\
T9 & Sí & Sí \\
T10 & Sí & Sí \\
\bottomrule
\end{tabular}
\end{mitabla}

\textit{Paso 2. Tabla de contingencia} (frecuencia observada $FO_{ij}$, $N = 148$). Cada celda cuenta a los titulados que respondieron a la vez la fila ($X$) y la columna ($Y$).

\begin{mitabla}{C11: tabla de contingencia, frecuencia observada $FO_{ij}$}
\begin{tabular}{L{3.8cm}C{3.00cm}C{3.00cm}|c}
\toprule
$X$: Formación complementaria \textbackslash{} $Y$: Interés en posgrado & Sí & No & Total \\
\midrule
Sí & 86 & 12 & 98 \\
No & 37 & 13 & 50 \\
\midrule
Total & 123 & 25 & $N = 148$ \\
\bottomrule
\end{tabular}
\end{mitabla}

\textbf{Mapa de calor.} La misma tabla de contingencia, pintada por color en vez de por número:

\begin{mitabla}{C11: mapa de calor de $FO_{ij}$}
\begin{tabular}{L{3.8cm}C{3.00cm}C{3.00cm}}
\toprule
$X$: Formación complementaria \textbackslash{} $Y$: Interés en posgrado & Sí & No \\
\midrule
Sí & \cellcolor{gray!75}86 & \cellcolor{gray!10}12 \\
No & \cellcolor{gray!32}37 & \cellcolor{gray!11}13 \\
\bottomrule
\end{tabular}
\end{mitabla}

\textit{Paso 3. Frecuencia esperada.} Es lo que habría en cada celda si $X$ e $Y$ fueran independientes. Por ejemplo, para la celda Sí y Sí:
\[
FE_{11} = \frac{(98)(123)}{148} = 81.45
\]

\begin{mitabla}{C11: frecuencia esperada $FE_{ij}$}
\begin{tabular}{L{3.8cm}C{3.00cm}C{3.00cm}}
\toprule
$X$: Formación complementaria \textbackslash{} $Y$: Interés en posgrado & Sí & No \\
\midrule
Sí & $98\cdot123/148 = 81.45$ & $98\cdot25/148 = 16.55$ \\
No & $50\cdot123/148 = 41.55$ & $50\cdot25/148 = 8.45$ \\
\bottomrule
\end{tabular}
\end{mitabla}

\textit{Paso 4. Estadístico chi-cuadrado.} En cada celda se calcula $\frac{(FO_{ij}-FE_{ij})^2}{FE_{ij}}$ y luego se suman todas.

\begin{mitabla}{C11: cálculo de $X_0^2$ celda por celda}
\begin{tabular}{L{3.2cm}C{3.20cm}C{3.20cm}|c}
\toprule
$X$: Formación complementaria \textbackslash{} $Y$: Interés en posgrado & Sí & No & Suma fila \\
\midrule
Sí & $\dfrac{(86-81.45)^2}{81.45}$\newline $= 0.2546$ & $\dfrac{(12-16.55)^2}{16.55}$\newline $= 1.2528$ & $1.5075$ \\[6pt]
No & $\dfrac{(37-41.55)^2}{41.55}$\newline $= 0.4991$ & $\dfrac{(13-8.45)^2}{8.45}$\newline $= 2.4555$ & $2.9546$ \\[6pt]
\bottomrule
\end{tabular}
\end{mitabla}

\[
X_0^2 = 1.5075 + 2.9546 = 4.4621
\]

\textit{Paso 5. Estadístico teórico.} Con $f = 2$ filas y $c = 2$ columnas: $gl = (2-1)(2-1) = 1$. De la tabla chi-cuadrada, con $\alpha = 0.05$: $X^2_{1,\,0.05} = 3.841$.

\textit{Paso 6. Relación entre estadísticos y conclusión.}

\begin{center}
\begin{tabular}{l}
$4.4621 \ge 3.841 \rightarrow$ falso\\
Si ($X_0^2 < X^2_{(f-1)(c-1),\,\alpha}$)\\
\quad ``Independencia''\\
Sino\\
\quad \fbox{NO ``Independencia''}\\
FinSi
\end{tabular}
\end{center}

Se rechaza $H_0$ y se acepta $H_1$: «Interés en posgrado» depende de «Formación complementaria». Porcentaje de cada fila en la columna ``Sí'': Sí $86/98 = 87.8\%$; No $37/50 = 74.0\%$.

\textit{Recomendación.} Quienes ya cursaron formación complementaria quieren seguir estudiando: ofrecer rutas de continuidad (diplomado, especialidad, maestría) a quienes ya terminaron un programa.

